\section{Finite arithmetic estimates}\label{app:arithmetic}
This section records the finite terms evaluated in Proposition~\ref{prop:arithmetic}. Research notes 02 and 03 give the local basis constructions and norm-product estimates; note 06 defines the column debit, effective slope, and refunds used to combine them. All sums below are finite. A maximum need not be found exactly: evaluating any admissible choice gives a valid lower bound.

\subsection{The retained contribution of a small inert prime}
Fix a prime $p<K$. Write $T(N,m)$ for the floor sum in \eqref{eq:T-floor}, and retain $W,M,W_2$ from Section~\ref{sec:comparison}. For an integer cut $0\leq L\leq S$, put
\[
 c_{K,p}=\frac{K+p-1}{2}-(p^2-1),\qquad
 q_L=\frac{(K-p)L(L+1)}{2M}.
\]
The fraction $q_L$ is the proportion of rows on which the stronger, but more expensive, derivative estimate is used. Define
\begin{align*}
 A_p(L)&=\frac{2(K-p)}{p^2}
       \sum_{t=1}^{L}t\bigl(K(S-t)+c_{K,p}\bigr),\\
 B_p(L)&=\frac{2(K-p)}{p^3}
       \sum_{t=L+1}^{S}t\bigl(K(S-t)+c_{K,p}\bigr),\\
 E_p(L)&=A_p(L)+B_p(L)-\frac{p^2}{2}(S-L)(K-p)\\
 &\hspace{1.4cm}+2(K-p)\left(T(L,p)-\sum_{a\geq1}T(L,p^{2a})\right),\\
 C_p(L)&=3\left(1+\frac1S\right)q_L
       \left(1-\frac{q_L\theta}{2}\right)\frac{p-1}{p}\frac JR
       +\boldsymbol{1}_{L>0}\frac{n\floor{p^2/4}}{pW}.
\end{align*}
With the notation of the local arithmetic notes (the formula following Lemma 4.3), an admissible retained row weight is
\begin{equation}\label{eq:wkeep}
 w_{\rm keep}(p)=
 \frac{W_2+2K\sum_{a\geq1}T(S,p^{2a})}{W}
 +\max_{0\leq L\leq S}
       \left(\frac{\pos{E_p(L)}}W-C_p(L)\right).
\end{equation}
It contributes $w_{\rm keep}(p)\log p$ after the common first-power column debit has been removed. The formula also applies to $p=2$ in the maximal order. The last term of $C_p$ uses $n=M/(\theta J)$ and is the same finite remainder as in the stated formula.

Here is the counting behind the two corrections. In the right part of the triangle, the rows can be grouped into strips of lengths $t$. Summing the extra Frobenius factors over a strip gives $A_p$ or $B_p$, according to which derivative estimate is used. Hermite's floor identity gives the loss $p^2(S-L)(K-p)/2$ at the incomplete ends of these strips. The derivative cost is controlled by the monotone function $j-\floor{j/p}$. Its exact prefix sum is
\begin{equation}\label{eq:small-column-prefix}
 \sum_{j=0}^{B-1}\left(j-\floor{j/p}\right)
 =\frac{p-1}{2p}B^2-\frac{r_B(p-r_B)}{2p},
 \qquad r_B=B\bmod p.
\end{equation}
Each derivative order is available $n$ times. Filling the largest orders first and using $r_B(p-r_B)\leq\floor{p^2/4}$ gives $C_p(L)$. Thus the factor $1+1/S$ and the final remainder in that expression are finite counting terms; omitting them would strengthen the bound without justification.

\subsection{Prime powers and the dependence on the selected columns}
For a prime-power cutoff $Q$, put
\[
 \sigma_2(Q)=\sum_{\substack{p\text{ prime},\ a\geq2\\p^a\leq Q}}
                       \frac{\log p}{p^a},\qquad
 H_Q(j)=\sum_{\substack{p\text{ prime},\ a\geq2\\p^a\leq Q}}
                      \left\{\frac j{p^a}\right\}\log p.
\]
Legendre's formula shows that these are respectively the linear part and the remainder of the higher-prime-power contribution. On a finite comparison, $Q=J-1$ retains every nonzero floor. The sorted capacity formula \eqref{eq:sorted-remainder} then bounds the contribution of $H_Q$ exactly. The non-inert-prime refunds in the hybrid arithmetic estimate modify both the slope and the remainder; their interval bounds are included before applying the same capacity argument.

For the unbounded range, we fix $Q=1000$. Discarding further factorial savings is legitimate because they have the favorable sign. There is a convenient bound that avoids sorting the joint remainder. For each $q=p^a\leq Q$, set
\[
 b_q=n\ceil{J/q},\qquad
 t_q=\min\left(q,\floor{M/b_q}\right),\qquad r_q=M-b_qt_q.
\]
Each residue modulo $q$ occurs among the columns at most $b_q$ times. Filling residues $q-1,q-2,\ldots$ in that order gives the normalized remainder bound
\[
 E_q=\frac1{Mq}\left[b_q\left(t_q(q-1)-\binom{t_q}{2}\right)
                    +r_q\max(0,q-1-t_q)\right].
\]
Consequently
\begin{equation}\label{eq:power-capacity}
 \frac2M\sum_C H_Q(j)\leq2\sum_{q=p^a\leq Q}E_q\log p.
\end{equation}
Compatibility between the residues at different prime powers is ignored here. The right side is therefore an upper bound, which is the direction required in \eqref{eq:column-remainder}.

\subsection{Norm values and removal of composite values}
The remaining arithmetic input is the explicit norm estimate in the local arithmetic notes, Section 6.4(3)--(4), including its two- and three-factor corrections. Its basic weight is $S_r(R,K,x)\log x/(KW)$, with $S_r$ defined in \eqref{eq:Sr}. On each interval where $\floor{R/x}$ is constant, this is a quadratic polynomial times $\log x$; the weight decreases for $x\geq3$. The norm-form count therefore gives an upper bound for the loss from possible non-inert primes by evaluating at lower norm endpoints.

For clarity, the multiplication used to certify composite values is ordinary multiplication in the quadratic order. If $\omega=(1+\sqrt{-D})/2$ and $m_D=(D+1)/4$, then
\[
 (x+y\omega)(x'+y'\omega)
 =(xx'-m_Dyy')+(xy'+x'y+yy')\omega.
\]
Products of nonunits have composite norm. Here an element is primitive if
it is not divisible in $\OO_D$ by any rational integer greater than one. In our range $D>4$, the units are $\{\pm1\}$. For a fixed norm $N=\prod_{i=1}^r p_i^{e_i}$, where the $p_i$ are distinct split rational primes and $e_i\geq1$, unique factorization gives exactly $2^{r-1}$ primitive elements modulo sign and conjugation. Indeed, for each $p_i=\pi_i\bar\pi_i$, primitivity forces a choice between $\pi_i^{e_i}$ and $\bar\pi_i^{e_i}$; conjugation pairs the resulting $2^r$ choices. The multiplicity depends on the number of distinct primes, not on $\sum_i e_i$. The norm-product estimate controls these multiplicities and the possible ordering changes of norm endpoints across a discriminant interval. These details are necessary: simply subtracting every product pair from a norm-value count would overcount. The numerical comparisons use the product count with these multiplicity bounds.

\section{Finite evaluation of the spectral estimate}\label{app:metric}
We give the bounds used to pass from the exact node estimates in Section~\ref{sec:metric} to a finite computation. The derivations are in research note 05, Section 5; note 06 gives the normalization used in the paper. They account separately for the signed row sum, the choice of derivative orders, and the number of columns selected at a node. At the finite comparisons, $R$ and $S$ are fixed, so the lower size parameters of the uniform estimates may be taken to equal $R$ and $S$.

\subsection{The signed row sum}
Write $L(x)=x\log x$, with $L(0)=0$, and define
\[
 f(x)=L(x)-L(a+x)+L(a)-x\log v,
 \qquad C_v=-a\log(1-v).
\]
For $aR>1$, integral comparison for the factorial and rising factorial gives
\[
 \lambda_{r,s}\leq Rf(r/R)+\beta s+\delta_R,
 \qquad
 \delta_R=\log(R+1)+1+\log\frac{a+1}{a-1/R}.
\]
In our application, $f(1)<0$. Since $f$ is convex, $f(0)=0$, and its minimum is $-C_v$, one has $f\leq0$ on $[0,1]$ and
$\operatorname{TV}((1-x)f(x))\leq3C_v$. The exact row multiplicity is
$S-\floor{r/K}\geq S(1-r/R)$. Multiplying by the negative function $f$ reverses this inequality, giving the upper bound
\begin{equation}\label{eq:row-finite}
 U_{\rm row}\leq
 3\int_0^1(1-x)f(x)\,\mathrm dx+\frac{9C_v}{R}
 +\frac\beta{2K}(1-S^{-2})
 +\frac{3(1+S^{-1})}{2R}\delta_R.
\end{equation}
The integral is elementary. The total-variation error bounds the left Riemann sum, while the last term accounts for the factorial errors at every row. Keeping this signed expression separate prevents an upper estimate for a number of rows from being multiplied by a negative cost.

\subsection{Column capacities and the forward bound}
Set $\varepsilon=J/R$. For each upper node radius $q>0$ and column scale $\tau$, define
\[
 g_\tau(z)=L(z)+L(a+z)-L(a)
          +2z(\ell-\log\tau-1),\qquad
 E_{\rm col}=\frac{\log(R+1)+1}{R}.
\]
The exact column price in \eqref{eq:column-prices} satisfies
$\nu_j/R\leq g_\tau(j/R)+E_{\rm col}$. Divide $[0,\varepsilon]$ into $N_z$ equal bins. Convexity of $g_\tau$ bounds its value on each bin by the larger endpoint value. Replace a negative upper value by zero, and give each bin capacity
\[
 c_z=\frac{\varepsilon}{N_z}+\frac1R.
\]
Let $C_{q,\tau}(t)$ be the greatest cost of selecting total mass $t$ subject to these bin capacities. It is computed by filling bins in decreasing order of their upper costs. Every selection of $m$ actual columns is feasible at $t=m/R$, so this gives an upper bound even when some exact prices are negative.

With $\delta_q=(\log q-\beta)_+$, $v_q=ve^{-\delta_q/K}$ and $w=v_qq$, the forward singular-product estimate divided by $R^2$ is at most
\begin{equation}\label{eq:forward-finite}
 C_{q,\tau}(t)+d_{q,\xi}t+t^2\log\xi,
 \qquad
 d_{q,\xi}=a\log\frac\Gamma{1-w}+\frac{\delta_q}{K}
 +\frac{\log R}{R}+E_{\rm col}-\frac{\log\xi}{R}.
\end{equation}
Here $\Gamma,\tau$ are obtained from $\xi$ as in Section~\ref{sec:metric}. The last correction retains the difference between $m(m-1)$ and $m^2$.

We must bound all possible occupancies, not just values on a grid. Divide $[0,\varepsilon]$ into $N_m$ cells. On a cell $[l,h]$, subtracting $a_0t$ from \eqref{eq:forward-finite} gives the whole-cell upper bound
\begin{equation}\label{eq:forward-cell}
 C_{q,\tau}(h)+
 \max\{(d_{q,\xi}-a_0)l,(d_{q,\xi}-a_0)h\}
 +l^2\log\xi.
\end{equation}
The first term uses monotonicity of the capacity cost; the last uses $\log\xi<0$. Thus no interpolation between sampled occupancies is assumed.

\subsection{The complementary bound}
For a scale $\kappa>0$, choose
\[
 d=\frac{\kappa}{1+\sqrt w},\qquad c_0=d\sqrt w,
 \qquad \tau=\frac{d(1-w)}{\sqrt{v_q}}.
\]
Use $G_0,s_0,r_0$ from Lemma~\ref{lem:complementary}, and set
\begin{align*}
 c_1&=-a\log G_0-\varepsilon
       (2\log r_0-\log s_0-|\log s_0|),\\
 c_2&=\log r_0-|\log s_0|<0,\\
 e_q&=-a\log(1-w)+\delta_q/K+\log R/R+E_{\rm col}
                       +\max(0,-\log d)/R.
\end{align*}
The complementary cost divided by $R^2$ is bounded above by
\begin{equation}\label{eq:complement-finite}
 C_{q,\tau}(t)+\varepsilon(a+\varepsilon)\log G_0
                   +(c_1+e_q)t+c_2t^2.
\end{equation}
This is obtained by expanding \eqref{eq:complementary-cost} and keeping the finite term $-\{t\log d+(\varepsilon-t)\log G_0\}/R$. Its unfavorable part is at most $t\max(0,-\log d)/R$. A whole-cell bound after subtracting $a_0t$ is obtained as in \eqref{eq:forward-cell}: use $C_{q,\tau}(h)$, the appropriate endpoint for the linear term, and $c_2l^2$.

\subsection{The Schur bound}
The Schur sum has a particularly convenient finite approximation. Put
\[
 \zeta=\ell+\tfrac12\log v_q-\log(1-w),\qquad
 z_*=a^{-2}e^{-2\zeta},\qquad
 I(x)=\tfrac12x^2\log x-\tfrac14x^2,
\]
Here $I(0)=0$ and $z_*$ agrees with the parameter in \eqref{eq:schur}:
$bc_*=ae^\zeta$. Define
\begin{align}
 g_\varepsilon(t)={}&t\bigl(\delta_q/K-a\log(1-w)\bigr)
 +(2\varepsilon t-t^2)(\zeta-1)+I(t)\notag\\
 &+I(a+2\varepsilon-t)-I(a+2\varepsilon-2t)
       -ta\log a-t^2\log(1-z_*).\label{eq:schur-finite}
\end{align}
Under $R>3$, $0<a\leq1$, $0<\varepsilon\leq1$,
$\varepsilon(a+2\varepsilon)\leq1$, and $ae^\zeta>1$, the logarithm of the right side of \eqref{eq:schur}, divided by $R^2$, is at most
\begin{equation}\label{eq:schur-error}
 g_\varepsilon(t)+E_0,\qquad
 E_0=\frac{\varepsilon}{R}
       (2\log R+2-\log\varepsilon-\log a).
\end{equation}
This is the spectral notes (Proposition~6). Its proof uses integral bounds for $\log i!$ and $\log(b)_j$; all powers of $\log R$ in the main term cancel. The displayed error includes the remaining factorial, mass, and Riemann-sum errors.

For a cell with midpoint $t_0$ and half-width $h_0$, the computation uses a tangent bound. Indeed,
\[
 g_\varepsilon''(t)=
 \log\frac{t(a+2\varepsilon-t)}
 {e^{2\zeta}(a+2\varepsilon-2t)^4(1-z_*)^2}.
\]
At each radial endpoint it verifies
$\varepsilon(a+\varepsilon)z_*<a^2(1-z_*)^2$, which implies concavity on $[0,\varepsilon]$. Consequently a whole-cell bound for $g_\varepsilon(t)-a_0t+E_0$ is
\begin{equation}\label{eq:schur-cell}
 g_\varepsilon(t_0)-a_0t_0
       +h_0|g_\varepsilon'(t_0)-a_0|+E_0.
\end{equation}
All values and slopes in this expression are enclosed with outward rounding.

\subsection{The finite parameter choices}
At each radius and occupancy cell, take the minimum of the whole-cell bounds just described, and then take the maximum over cells and zero. This is an admissible $\Phi(q)$ in \eqref{eq:H-occupancy}. The order of these operations matters: a bound is established on the entire cell before a spectral parameter or branch is selected. Finally insert the resulting $\Phi$ into \eqref{eq:metric-mesh}.

In addition to \eqref{eq:spectral-parameters}, the finite calculations use
\[
 \kappa\in\{5/2,3,4,6,10\},
\]
\begin{multline*}
 \xi,r_0\in\{1/100,3/100,3/50,1/10,3/20,1/5,3/10,2/5,\\
                     1/2,3/5,7/10,4/5,9/10,19/20,99/100,9999/10000\}.
\end{multline*}
Only positive upper radial endpoints are used; the exact node formulas at $q=0$ follow by continuity. The implementations also enlarge the derivative slope $\ell$ to a containing rational interval before evaluating the spectral bound. This is safe because each fixed selected minor is multiplied by $e^{2\ell\sum j}$ and hence increases with $\ell$. The actual interval endpoints and rounding records are retained with the numerical output.

\section{The unbounded range and numerical records}\label{app:global}
The comparison for $D\geq875$ uses the radial estimate in research note 04, the resultant rank lemma in note 01, and the local arithmetic estimates in note 02. It is useful to record why the numerical proof covers an unbounded range. A positive evaluation at one large value of $D$ would not suffice. What is needed, and what the calculation establishes, is a lower bound for the gap that increases with $\log D$.

\subsection{A uniform comparison above 5,760,000}
Set
\[
 D_0=5\,760\,000,\quad X=(D-3)/4,\quad X_0=(D_0-3)/4,
 \quad K=31,\quad c_R=23/10,\quad\varepsilon=29/500,
\]
and choose
\[
 S=\floor{c_R X/K},\quad R=KS,\quad J=\floor{\varepsilon R}.
\]
Write $S_0=106\,838$, $R_0=3\,311\,978$, $J_0=192\,094$ for their anchor lower bounds, and put $J_* =J_0+1$ and $M_0=R_0(S_0+1)/2$. The node number is rounded upward to a multiple of $800$ so that the resultant criterion
\[
 n>R(S+M_y)/J
\]
holds; here $M_y$ is an upper bound for the number of distinct imaginary coordinates of the nodes. The additional hypothesis
$\max(J,M_x)\leq R+KM_y$ of the local arithmetic and global comparison notes (Lemma~3.3) is also retained. The disk estimate gives $M_x\leq\floor{2\rho}+1$ and $M_y\leq2\floor{2\rho/\sqrt D}+1$.

To obtain a bound for the whole ray, we first bound $n/X$, the occupancy,
and the error in the disk radius by constants depending only on $D_0$.
The remaining dependence on $D$ will then occur through $\log(X/X_0)$. Put
\begin{align*}
 c_{R,-}&=c_R-K/X_0,&\varepsilon_-&=\varepsilon-1/R_0,\\
 \varepsilon_\rho&=\frac{4\sqrt{1.02K/\varepsilon_-}}
                {\sqrt{2\pi S_0\sqrt{D_0}}}
       +\frac{1+\sqrt{1+4/D_0}}{S_0},&
 \gamma&=1+\varepsilon_\rho+800\varepsilon/S_0.
\end{align*}
One checks $\gamma<1.002779878983<1.02$. Rounding
$(1+\varepsilon_\rho)RS/J$ strictly upward to a multiple of $800$ gives an admissible $n$ with
$n\leq\gamma RS/J$. This closes the preliminary radius estimate made with $n\leq1.02RS/J$. The other resultant hypothesis follows from $J<R$ and $2\rho+1<R$.
Indeed, the radius bounds below give
\[
 \frac{2\rho+1}{R}\leq
 \frac{2e^{C_+}}{c_{R,-}}X^{-1/4}+\frac1{c_{R,-}X}.
\]
The right side decreases with $X$ and is less than one at $X_0$.

Take the rational lower occupancy bound
$\theta_-=398891129/800000000$ and put $\theta_+=(1+1/S_0)/2$. Define
\begin{align*}
 \nu_-&=c_{R,-}/(K\varepsilon),&
 \nu_+&=\gamma c_R/(K\varepsilon_-),\\
 \eta_+&=\frac{\sqrt{D_0+4}/4}
               {\sqrt{\nu_-X_0\sqrt{D_0}/(2\pi)}},&
 C_-&=\tfrac12\log(\nu_-/\pi),\\
 C_+&=\tfrac12\log\bigl(\nu_+\sqrt{1+3/(4X_0)}/\pi\bigr)
                           +\log(1+\eta_+).
\end{align*}
Then $\theta_-\leq M/(nJ)\leq\theta_+$, $\nu_-X\leq n\leq\nu_+X$, and
\begin{equation}\label{eq:tail-radius}
 \tfrac34\log X+C_-\leq\log\rho\leq\tfrac34\log X+C_+.
\end{equation}
In particular $\eta_+<0.022619736410$. The factor $\sqrt{D_0}$ in its definition comes from the lattice covolume and is essential.

\begin{proposition}[Uniform local comparison]
\label{prop:uniform-tail}
For these choices, the local arithmetic estimates (Sections 4.3 and 6.1 of the research notes) and the radial estimate (Section 5.2) imply
\begin{align}
 \frac{L_{\rm arith}(D)}{W}&\geq A_0+\log(X/X_0),\label{eq:tail-arith}\\
 \frac{U_{\rm arch}(D)}{W}&\leq U_0+\tfrac34(A+B)\log(X/X_0),\label{eq:tail-metric}
\end{align}
where
\begin{align*}
 A&=1+\frac1K+\frac{3(K-1)}{2KS_0}+\frac{K-5}{2KS_0^2},\\
 B&=\frac{3(S_0+1)J_*}{R_0S_0}(1-\theta_-/2)
                  +\frac{3\theta_+J_*}{R_0^2S_0},
\end{align*}
and the directed evaluations give
\[
 A_0>10.8131594548034104832,
 \qquad U_0<10.8128420196002531239.
\]
These bounds apply to fundamental integer discriminants with class number one throughout $D\geq D_0$.
\end{proposition}

We next explain the constants in this estimate and their uniformity in $D$. This also identifies the part of the global argument that is distinct from the spectral estimate used near $499$. With $\sigma=\sigma_2(1000)$, put
\[
 \ell_- =\tfrac14\log X_0+\log c_{R,-} -C_+-\sigma,
 \qquad
 \ell_+ =\tfrac14\log X_0+\log c_R-C_- -\sigma.
\]
The actual derivative slope is bounded between these endpoints plus
$\tfrac14\log(X/X_0)$. The upper constant is evaluated as
\begin{align}
 U_0={}& A(\tfrac34\log X_0+C_+)+B\ell_+
       -F_*-S_{\rm der}\notag\\
 &+\frac{3(S_0+1)}{2R_0S_0}c_j
       +\frac{3(S_0+1)}{R_0S_0}\log M_0.\label{eq:tail-upper-constant}
\end{align}
The terms $c_j$, $S_{\rm der}$ and $F_*$ respectively bound the prime-power remainder, the falling-factorial saving, and the coupled radial and derivative saving. The evaluated bounds used here are
\begin{center}
\begin{tabular}{lr}
\toprule
Quantity & Directed bound\\
\midrule
$c_j$ & $<60.016864219225$\\
$S_{\rm der}$ & $>0.1803971924811329636$\\
$F_*$ & $>0.0825081841978627281$\\
\bottomrule
\end{tabular}
\end{center}
The formula \eqref{eq:power-capacity}, with the uniform occupancy and size bounds, gives $c_j$. The falling-factorial saving is the explicit lower Riemann sum with $3200$ intervals in Section 5.1 of the research notes. Its integrand is the improvement over the linear bound obtained from
\[
 \log\frac{\fall rj}{R^j}
 \leq R\left[x\log x-(x-y)\log(x-y)-y\right],
 \qquad x=\min(1,(r+1)/R),\quad y=j/R.
\]
The quantile bounds on the triangle fix a lower $y$ and an upper $x$ on each interval. The saving increases with the lower derivative slope and decreases with the upper row degree. Evaluating these at their anchor bounds therefore gives a uniform saving.

For $F_*$, Lemma 5.2 of the research notes gives a convex function of $800$ node occupancies, each in $[0,1]$, with their sum prescribed. It consists of a convex radial term in cumulative occupancies, positive quadratic terms, and positive squared ramps. The supplied rational vector has denominator $10^6$ and coordinate sum $398891129/10^6=800\theta_-$. Evaluating the function and a supporting hyperplane at this vector gives
\[
 F(w_*)>0.08250828755716835749,\qquad
 F(w_*)-\inf F<0.00000010335930562937,
\]
which implies the stated lower bound for $F_*$. The linear minimization in the supporting-hyperplane calculation is exact: fill the coordinates with the smallest gradient coefficients first, subject to the box and sum constraints. This gives a lower bound for the convex function from a single rational
vector. The defining formula is Lemma 5.2 of the research notes; the vector and
its supporting hyperplane are included with the supplementary calculations.

On the arithmetic side, all primes retained in this tail are below the guaranteed inert cutoff $X$. The local row weights are combined with the explicit prime estimates of Rosser--Schoenfeld \cite[Theorems~4, 6 and 18]{rosserschoenfeld}. In particular, writing
\[
 E_M=-\gamma_{\rm Euler}-\sum_p\frac{\log p}{p(p-1)},
\]
one uses
\[
 \sum_{p\leq X}\frac{\log p}{p}
 \geq\log X+E_M-\frac1{2\log X}.
\]
The constant is enclosed in
$[-1.33258227573322089,-1.33258227573322087]$.
The small-prime weights and the even-power tower are evaluated from their finite lower formulas at $S_0$; the tail estimates in Section 6.1 of the research notes bound each retained contribution
from below by its value at $S_0$. The negative prime-sum errors and floor losses decrease as $X$ increases. This gives \eqref{eq:tail-arith}, with the coefficient of $\log(X/X_0)$ equal to one.

The uniform bounds for the finite remainders in
\eqref{eq:tail-upper-constant} are justified in Section 6.1 of the research notes; the
underlying terms decrease with $S$. The geometry grows by at most $\tfrac34A\log(X/X_0)$, and the derivative term by at most $\tfrac14B\log(X/X_0)$. The implementation conservatively charges $\tfrac34B$ instead of $\tfrac14B$, yielding \eqref{eq:tail-metric}. Consequently
\begin{equation}\label{eq:tail-positive}
 \frac{L_{\rm arith}-U_{\rm arch}}W
 >0.00031743+0.127829\log(X/X_0)>0.
\end{equation}
This contradicts the nonzero determinant comparison on the entire ray. The proof of Proposition~\ref{prop:uniform-tail} uses the local rank, divisibility and radial lemmas; the constants above and those uniformity estimates give the displayed
logarithmic lower bound.

\subsection{Finite coverage}
Below $D_0$, fixed parameter choices cover $[875,D_0]$ by $101$ adjoining closed intervals. The three stages and their minimum gaps are listed in Table~\ref{tab:coverage}.
The first stage uses norm-form bounds for non-inert primes, the second exact finite prime sums, and the third the explicit prime estimates of \cite{rosserschoenfeld}. Parameters are fixed on each interval; the radius and slope are enclosed over the whole interval. The upper endpoint of one interval equals the lower endpoint of the next. The complete schedule, including all integer parameters and directed gap endpoints, is supplied as \code{verification/results/global875-schedule.csv}. Each comparison therefore establishes a positive gap throughout its interval.

The supplementary directory \code{verification/} contains the programs and
numerical results for these comparisons, the fifteen individual values above
$499$, and the calculation on $[499,503]$. Its short reproduction guide gives
the commands for each calculation. The arithmetic modules used at $499$ and at the preceding comparisons are included. The decimal and binary evaluations share some intermediate
routines. They provide numerical enclosures for the estimates above; the
mathematical inputs are the lemmas stated or cited in the paper.
